\section{Выводы и обсуждение результатов}

В ходе выполнения лабораторной работы достигнуты основные цели исследования:  
\begin{itemize}  
    \item Изучены явления дифракции Френеля и Фраунгофера на одной щели. Экспериментально подтверждены ключевые закономерности:  
    \begin{itemize}  
        \item Для дифракции Френеля измеренная ширина щели \( b = 0.296 \pm 0.024 \, \text{мм} \) совпала с прямым измерением (\( 0.3 \, \text{мм} \), \( \varepsilon = 1\% \)).  
        \item Для дифракции Фраунгофера получена ширина щели \( b = 0.353 \pm 0.007 \, \text{мм} \), что согласуется с измеренным значением \( 0.35 \, \text{мм} \) (\( \varepsilon = 0.86\% \)).  
    \end{itemize}  

    \item Проверены теоретические соотношения для положения минимумов. Линейные зависимости \( 4\lambda m(1/z) \) и \( x_m(m) \) подтверждены с высокой точностью (\( R^2 = 0.998 \)).  

    \item Исследовано влияние дифракции на разрешающую способность. Теоретический предел \( b_{\text{крит}} = 0.092 \, \text{мм} \), однако экспериментальное исчезновение полос при \( b_0 = 1.31 \, \text{мм} \) указывает на неточность измерения входной щели.
\end{itemize}  

Работа продемонстрировала фундаментальную связь между геометрией оптической системы и дифракционными эффектами. Все основные теоретические предсказания подтверждены в пределах погрешностей измерений. Расхождения в исследовании разрешающей способности требуют дополнительного анализа условий эксперимента.  

\subsection*{Основные результаты}  
\begin{itemize}  
    \item Дифракция Френеля: \( b = 0.296 \pm 0.024 \, \text{мм} \), \( \chi^2 = 2.1 \) (\( \chi^2_{\text{крит}} = 7.815 \)).  
    \item Дифракция Фраунгофера: \( b = 0.353 \pm 0.007 \, \text{мм} \), \( R^2 = 0.998 \).  
    \item Разрешающая способность: \( b_0 = 1.31 \, \text{мм} \), \( b_{\text{крит}}^{\text{теор}} = 0.092 \, \text{мм} \).  
\end{itemize}  
